%! AP Statistics — Unit 1, Lesson 1.4 — Guided Notes ANSWER KEY
\documentclass[10pt]{article}
\usepackage{apstats-article}
\usepackage{apstats-key}

%=============================================================================
\begin{document}

\pageheader{Unit 1, Lesson 1.4}{Guided Notes --- Answer Key}
\begin{tabularx}{\linewidth}{X X X}
Name:\;\ans{Key} & Date:\; & Period:\;
\end{tabularx}

\vspace{0.12in}

\begin{objectivebox}
\small By the end of this lesson, I will be able to\dots\\[4pt]
\ansline{construct a correctly labeled bar chart (frequency and relative frequency) from a table.}\\
\ansline{construct a pie chart by computing sector angles as relative frequency $\times\,360^\circ$.}\\
\ansline{identify misleading features in a graph and explain how to fix them.}
\end{objectivebox}

\vspace{0.1in}

\begin{vocabbox}
\small
\noindent\textbf{\textcolor{navy}{Bar chart (bar graph):}}\\
\ans{A graph with separated bars whose heights represent the frequency or relative frequency of each category of a categorical variable}\\[6pt]

\noindent\textbf{\textcolor{navy}{Relative frequency bar chart:}}\\
\ans{A bar chart with relative frequency (proportion or percent) on the vertical axis; preferred when comparing groups of different sizes}\\[6pt]

\noindent\textbf{\textcolor{navy}{Pie chart:}}\\
\ans{A circle divided into sectors, each representing one category; sector sizes are proportional to relative frequencies}\\[6pt]

\noindent\textbf{\textcolor{navy}{Segmented bar chart:}}\\
\ans{A bar chart with bars divided into colored segments to show the composition of each group; useful for comparing distributions across groups}\\[6pt]

\noindent\textbf{\textcolor{navy}{Misleading graph:}}\\
\ans{A visual display that distorts or exaggerates differences through a non-zero y-axis, 3-D effects, inconsistent scale, or other deceptive design choices}\\[6pt]
\end{vocabbox}

\vspace{0.1in}

\begin{hookbox}
\small\textbf{What's Wrong With These Graphs?}

\vspace{8pt}
\begin{enumerate}[leftmargin=1.5em, itemsep=6pt]
  \item \ans{No. The actual difference is $1000 - 970 = 30$ units, which is only
        3\% relative to the y-axis midpoint. Because the axis starts at 950, the
        bar for Product Y looks roughly twice as tall as Product X, which
        exaggerates the difference dramatically.}

  \begin{teachernote}
  Prompt: ``If the axis started at 0, would the bars look so different?'' Have
  students sketch both versions mentally. Product X: 970/1000 bar height vs.\
  Product Y: 1000/1000 --- barely distinguishable.
  \end{teachernote}

  \item \ans{Start the y-axis at 0 so bar heights accurately reflect the data
        values. Add a clear axis label with units.}

  \item \ans{The 3-D perspective makes the near slices appear larger than they
        are. A reader may overestimate the relative frequency of the categories
        at the front of the chart. Flat (2-D) pie charts avoid this distortion.}
\end{enumerate}
\end{hookbox}

\vspace{0.1in}

\begin{notesbox}{Constructing a Bar Chart --- Four Requirements}
\small

\begin{multicols}{2}

\textbf{\textcolor{navy}{A bar chart must have:}}
\begin{enumerate}[leftmargin=1.4em, itemsep=4pt]
  \item A \ans{title}.
  \item A labeled \ans{horizontal (x-)} axis listing the categories.
  \item A labeled \ans{vertical (y-)} axis with a scale starting at \ans{0}.
  \item \ans{Gaps} between bars.
\end{enumerate}

\vspace{4pt}
{\color{navy}\small\textit{Note:} The y-axis may show frequency \textit{or}
relative frequency --- label it clearly!}

\columnbreak

\textbf{\textcolor{navy}{Frequency vs.\ Relative Frequency}}

\vspace{4pt}
\begin{tabular}{@{}p{0.42\linewidth} p{0.42\linewidth}@{}}
\textbf{Frequency bar chart} & \textbf{Rel.\ freq.\ bar chart} \\
y-axis: \ans{counts} & y-axis: \ans{proportions/\%} \\[4pt]
Use when: \ans{one group} & Use when: \ans{comparing} \\
& \ans{groups} \\
\end{tabular}

\end{multicols}
\end{notesbox}

\vspace{0.1in}

\begin{notesbox}{Worked Example --- Pet Ownership Bar Chart ($n = 40$)}
\small

\noindent Data: Dog~30\%, Cat~20\%, Fish~15\%, Bird~10\%, None~25\%.

\vspace{6pt}
\begin{teachernote}
Correct bar chart: 5 separated bars labeled Dog, Cat, Fish, Bird, None on x-axis.
Y-axis labeled ``Relative Frequency'' (or ``Proportion''), scale 0 to 0.35
(or 0 to 35\%), tick marks at 0, 0.05 (5\%), 0.10, etc.
Title: e.g., ``Distribution of Pet Type for 40 Households.''
Bars: Dog 0.30, Cat 0.20, Fish 0.15, Bird 0.10, None 0.25.
Bars must NOT touch.
\end{teachernote}

\vspace{8pt}
\textbf{AP-style sentence:}\\[4pt]
\ans{Of the 40 households surveyed, dog was the most common pet type (30\%),
followed by ``none'' (25\%), cat (20\%), fish (15\%), and bird (10\% --- the
least common).}

\begin{teachernote}
Require: total $n$, variable in context, most common and least common categories
with relative frequencies. Deduct if student writes only counts.
\end{teachernote}
\end{notesbox}

\newpage

\begin{notesbox}{Constructing a Pie Chart}
\small
\begin{multicols}{2}

\textbf{\textcolor{navy}{Sector angle:}}
\[
  \text{Sector angle} = \text{relative frequency} \times \textcolor{keyred}{\mathbf{360^\circ}}
\]

\vspace{4pt}
\textbf{\textcolor{navy}{Pet data --- angles:}}

\vspace{4pt}
\renewcommand{\arraystretch}{1.6}
\begin{tabular}{@{} >{\bfseries}p{0.22\linewidth}
                     >{\centering\arraybackslash}p{0.18\linewidth}
                     >{\centering\arraybackslash}p{0.24\linewidth} @{}}
\rowcolor{sky}
\textbf{Category} & \textbf{Rel.\ Freq.} & \textbf{Angle (${}^\circ$)} \\
\midrule
Dog  & 0.30 & \ans{108} \\
Cat  & 0.20 & \ans{72} \\
Fish & 0.15 & \ans{54} \\
Bird & 0.10 & \ans{36} \\
None & 0.25 & \ans{90} \\
\midrule
\rowcolor{sky}\textbf{Total} & 1.00 & \ans{360} \\
\end{tabular}

\columnbreak

\textbf{\textcolor{navy}{When to use a pie chart:}}
\begin{itemize}[itemsep=3pt]
  \item Emphasizes \ans{part-to-whole} relationships.
  \item Works best with \ans{few} categories (6 or fewer).
  \item \textit{Not} ideal for \ans{many} categories.
\end{itemize}

\vspace{6pt}
\textbf{\textcolor{navy}{Bar vs.\ Pie:}}\\[4pt]
Bar charts are generally preferred because they are easier to \ans{compare} and do not rely on judging \ans{angle/area} sizes.

\end{multicols}
\end{notesbox}

\vspace{0.1in}

\begin{notesbox}{Misleading Graphs --- Three Common Tricks}
\small

\begin{multicols}{2}

\textbf{Trick 1: Non-zero y-axis}\\
The y-axis starts above \ans{0}.\\[4pt]
\textit{Fix:} \ans{Start the y-axis at 0.}

\vspace{8pt}
\textbf{Trick 2: 3-D or perspective effects}\\
Slices/bars near the viewer appear \ans{larger}.\\[4pt]
\textit{Fix:} \ans{Use a flat (2-D) display.}

\columnbreak

\textbf{Trick 3: Inconsistent scale}\\
\textit{Fix:} \ans{Use equal intervals on all axes; keep bar widths uniform.}

\vspace{8pt}
\textbf{\textcolor{navy}{AP-style answer template:}}\\[4pt]
``The graph is misleading because [specific feature]. This makes [category X]
appear [larger/smaller] than it actually is. To fix it, [specific correction].''

\end{multicols}
\end{notesbox}

\vspace{0.1in}

\begin{practicebox}
\small
\textbf{TV Genre Survey ($n = 80$):}

\vspace{8pt}
\renewcommand{\arraystretch}{1.8}
\begin{tabularx}{\linewidth}{@{} >{\bfseries}p{0.20\linewidth}
                                    C{0.18\linewidth}
                                    C{0.24\linewidth}
                                    C{0.22\linewidth} @{}}
\rowcolor{sky}
\textbf{Genre} & \textbf{Frequency} & \textbf{Rel.\ Frequency} & \textbf{Angle (${}^\circ$)} \\
\midrule
News    & 32 & \ans{0.40 (40\%)} & \ans{144} \\
Sports  & 20 & \ans{0.25 (25\%)} & \ans{90}  \\
Reality & 16 & \ans{0.20 (20\%)} & \ans{72}  \\
Other   & 12 & \ans{0.15 (15\%)} & \ans{54}  \\
\midrule
\rowcolor{sky}\textbf{Total} & \ans{80} & \ans{1.00} & \ans{360} \\
\end{tabularx}

\vspace{10pt}
\begin{enumerate}[leftmargin=1.5em, itemsep=6pt]
  \item \textit{[Sketch --- see Teacher Note.]}
  \begin{teachernote}
  Correct bar chart: 4 bars for News (tallest, 40\%), Sports (25\%), Reality (20\%),
  Other (shortest, 15\%). Y-axis labeled ``Relative Frequency'' from 0 to at least
  0.40. Bars do not touch. Title present. Accept either decimal or percent axis.
  \end{teachernote}

  \item \ans{Of the 80 adults surveyed, news was the most preferred TV genre
        (40\%), followed by sports (25\%), reality TV (20\%), and other (15\%).}

  \item \ans{Starting the y-axis at 10\% (0.10) instead of 0 makes the News bar
        appear about 4$\times$ taller than Sports instead of the true ratio of
        $40\%/25\% = 1.6\times$. This exaggerates the dominance of news viewing
        and is misleading.}
\end{enumerate}
\end{practicebox}

\vspace{0.1in}

\begin{spiralbox}
\small
\begin{multicols}{2}

\textbf{Connections:}
\begin{itemize}[itemsep=2pt]
  \item \textbf{Lesson 1.5:} Histograms have \textit{touching} bars for continuous data.
  \item \textbf{Unit 2:} Segmented bar charts for two categorical variables.
  \item \textbf{AP Exam:} Graph + written interpretation in free-response.
\end{itemize}

\columnbreak

\textbf{Big Idea:}\\[4pt]
\ans{Always check that the y-axis starts at 0. A non-zero axis is the single
most common way a bar chart lies. Also check for a title and labeled axes ---
without them, the graph cannot be interpreted reliably.}

\end{multicols}
\end{spiralbox}

\end{document}
